% !TeX program = lualatex
% Compile twice: lualatex 115.tex
% All references and prose are embedded; no BibTeX database or external inputs.
\documentclass[11pt,a4paper]{article}
\usepackage[margin=24mm,headheight=14pt]{geometry}
\usepackage{fontspec}
\setmainfont{Latin Modern Roman}
\setsansfont{Latin Modern Sans}
\setmonofont{Latin Modern Mono}
\usepackage{amsmath,amssymb,mathtools}
\usepackage{unicode-math}
\setmathfont{Latin Modern Math}
\usepackage{microtype}
\usepackage{enumitem,xurl,needspace}
\usepackage[dvipsnames]{xcolor}
\usepackage{hyperref}
\hypersetup{unicode=true,colorlinks=true,linkcolor=MidnightBlue,urlcolor=MidnightBlue,
 pdftitle={Research attempt: Problem 115},pdfauthor={Research notebook},bookmarksnumbered=true}
\usepackage{bookmark}
\usepackage{tocloft}
\setlength{\cftsecnumwidth}{3em}
\cftsetpnumwidth{2.5em}
\cftsetrmarg{4em}
\usepackage{titlesec}
\titleformat{\section}{\Large\bfseries\raggedright}{\thesection}{0.7em}{}
\usepackage{fancyhdr}
\pagestyle{fancy}\fancyhf{}
\fancyhead[L]{\small\sffamily Open Problems Math | Research notebook}
\fancyhead[R]{\small Problem 115 | 27 September 2026}
\fancyfoot[C]{\thepage}
\setlength{\parindent}{0pt}
\setlength{\parskip}{3pt plus 1pt}
\setlength{\emergencystretch}{3em}
\setcounter{tocdepth}{1}
\setcounter{secnumdepth}{1}
\setlist[itemize]{leftmargin=3em,itemsep=5pt,topsep=4pt,parsep=0pt}
\allowdisplaybreaks
\newcommand{\N}{\mathbb{N}}
\newcommand{\Z}{\mathbb{Z}}
\newcommand{\Q}{\mathbb{Q}}
\newcommand{\R}{\mathbb{R}}
\newcommand{\C}{\mathbb{C}}
\newcommand{\F}{\mathbb{F}}
\newcommand{\A}{\mathbb{A}}
\newcommand{\PP}{\mathbb P}
\newcommand{\E}{\mathbb{E}}
\newcommand{\eps}{\varepsilon}
\newcommand{\dd}{\,\mathrm d}
\newcommand{\sref}[2]{\hyperlink{source:#1:#2}{[S#2]}}
\newcommand{\eref}[2]{\hyperlink{extra:#1:#2}{[E#2]}}
% Turn the next switch off for a shorter reading copy. The quotations preserve
% every exactTarget field, including qualifications omitted from a short summary.
\newif\ifshowcatalogue
\showcataloguetrue
\newcommand{\cataloguescope}[1]{\ifshowcatalogue
 \paragraph{Exact scope in the supplied catalogue.}
 \begin{quote}\small #1\end{quote}\fi}
\usepackage{etoolbox}
\AtBeginEnvironment{itemize}{\small}
\numberwithin{equation}{section}
% Standalone export. Original section 115 preserved verbatim outside marked additions.
\begin{document}
\setcounter{section}{114}
% ================================================================
% problemId: problem.the-borel-conjecture-on-topological-rigidity-of-aspherical-manifolds
% source releaseRank: 115; source releaseStatus: open_with_solved_subcases
% exactTarget SHA-256: 852713e775532076c94e3b61386c6f567858cdb52567563f12098e31e0c0c8da
\clearpage
\section{\texorpdfstring{The Borel Conjecture on Topological Rigidity of Aspherical Manifolds}{The Borel Conjecture on Topological Rigidity of Aspherical Manifolds}}\label{115}
\subsection{Definitions and mathematical statement}
A connected manifold is aspherical if its universal cover is contractible, equivalently its higher homotopy groups vanish. Given closed aspherical manifolds $M,N$ and a homotopy equivalence $f:M\to N$, the conjecture asks for a homeomorphism $h:M\to N$ with
\[
 f\simeq h.
\]
The specific homotopy class must contain a homeomorphism; merely knowing that the two manifolds happen to be homeomorphic is weaker. The category is topological, so smooth rigidity or a diffeomorphism is not requested.
\par\smallskip\textit{Formulation and concept references:} \sref{115}{1}.
\cataloguescope{Prove that every homotopy equivalence between closed aspherical manifolds is homotopic to a homeomorphism.}
\subsection{Short English statement}
Does every homotopy equivalence between closed aspherical manifolds come from a homeomorphism up to homotopy?
% BEGIN RESEARCH ADDITION 115
\subsection{Research attempt}

\paragraph{Current frontier and exact scope.}
The source's Section~2 distinguishes rigidity of a given homotopy class
from mere classification of manifolds. It also distinguishes topological
rigidity from smooth rigidity: exotic smooth structures are not
counterexamples here. Bartels--L\"uck prove the requested conclusion in
dimension at least five for, among others, hyperbolic and finite-dimensional
$\operatorname{CAT}(0)$ groups. Their Proposition~0.3 states the precise
$K$- and $L$-theoretic hypotheses, including orientation twists; its
four-dimensional extension requires a good fundamental group in
Freedman's sense. The original target imposes none of these restrictions.
\sref{115}{1}\eref{115}{1}

A useful recent caution comes from Gu's August 2026 revision on
$\mathcal Z$-compactifications: removing an interval factor requires
additional compatibility, and applying compactifications to universal
covers requires further control. This is adjacent work, not a proof of
Borel rigidity. Its corrected hypotheses warn against a proposed shortcut
from contractibility to a controlled homeomorphism. \eref{115}{2}

\paragraph{Attack route.}
Three routes are tested. First, retain exactly the two integral assembly
properties that would kill the geometric obstruction. Second, try to
construct a homeomorphism on universal covers and descend it equivariantly.
Third, exploit the explicit linear realization of automorphisms for tori.
The first produces a conditional reduction; the second fails without
small-scale control; the third supplies a restricted verification, not a
classification of all manifolds homotopy equivalent to a torus.

\paragraph{Attempt: the two-degree obstruction calculation.}
Let $N$ be a closed aspherical topological $n$-manifold, $n\geq5$,
$\pi=\pi_1(N)$, and $w:\pi\to\{\pm1\}$ its orientation character.
Assume $\operatorname{Wh}(\pi)=0$. A homotopy equivalence to $N$ is then
simple, so it defines an element of the topological simple structure set
$\mathcal S^{\mathrm{TOP}}(N)$. Its basepoint is the identity of $N$;
being the basepoint means precisely that the homotopy class contains a
homeomorphism, not just that its source happens to be homeomorphic to $N$.
Use the orientation-twisted quadratic $L$-spectrum throughout.
The relevant part of the surgery sequence is
\begin{equation}\label{eq:115:surgery}
\begin{split}
 H_{n+1}(N;\mathbf L\langle1\rangle)&\xrightarrow{A^c_{n+1}}
 L^s_{n+1}(\Z\pi,w)\longrightarrow\mathcal S^{\mathrm{TOP}}(N)\\
 &\xrightarrow{\eta}H_n(N;\mathbf L\langle1\rangle)
 \xrightarrow{A^c_n}L^s_n(\Z\pi,w).
\end{split}
\end{equation}
Here $\mathbf L\langle1\rangle$ retains the positive homotopy groups
of the periodic spectrum and kills degrees at most zero. The middle
sequence is interpreted with the usual structure-set action, not as an
exact sequence of abelian groups. This is the geometric/algebraic surgery
input used in Bartels--L\"uck's Proposition~0.3. \eref{115}{1}

\textbf{Lemma 1 (a sufficient two-degree criterion).}
If $A^c_n$ is injective and $A^c_{n+1}$ is surjective, then
$\mathcal S^{\mathrm{TOP}}(N)$ has one element.

\emph{Proof.} For every structure $s$, exactness gives
$A^c_n\eta(s)=0$. Injectivity gives $\eta(s)=0$. Exactness at the
structure set says that every such $s$ lies in the orbit of the basepoint
under $L^s_{n+1}(\Z\pi,w)$. Elements coming from $A^c_{n+1}$ fix that
basepoint. Surjectivity therefore makes the entire orbit a singleton.
Every $s$ is the basepoint. $\square$

This identifies why an injectivity theorem alone is not enough for this
argument: it eliminates the normal invariant but does not by itself
eliminate the residual surgery action. It also identifies why rational
injectivity is weaker still: it need not kill integral torsion in the
normal invariant.

\paragraph{Attempt: check the connective cutoff rather than dropping it.}
Let $\mathbf C$ be the cofiber of
$\mathbf L\langle1\rangle\to\mathbf L$. Then
$\pi_j\mathbf C=0$ for $j>0$. The homological dimension of an
$n$-manifold, with the local coefficients occurring here, is at most $n$.
The generalized-homology spectral sequence
\[
 H_a(N;\pi_b\mathbf C)\Longrightarrow H_{a+b}(N;\mathbf C)
\]
therefore has no nonzero terms of total degree greater than $n$:
$a\leq n$ and $b\leq0$. The base dimension is bounded, so no unbounded
filtration or convergence interchange is being used. Consequently
\[
 H_{n+2}(N;\mathbf C)=H_{n+1}(N;\mathbf C)=0.
\]
The cofiber long exact sequence now gives
\begin{equation}\label{eq:115:comparison}
\begin{aligned}
 H_{n+1}(N;\mathbf L\langle1\rangle)&\xrightarrow{\ \simeq\ }
 H_{n+1}(N;\mathbf L),\\
 H_n(N;\mathbf L\langle1\rangle)&\hookrightarrow H_n(N;\mathbf L).
\end{aligned}
\end{equation}
Thus isomorphisms of the \emph{periodic simple} assembly maps in degrees
$n,n+1$ imply the hypotheses of Lemma~1. The degree-$n$ comparison need
not be surjective; claiming an isomorphism there would be unnecessary
and unjustified.

Asphericity supplies $N\simeq B\pi$. To obtain these periodic simple
assembly isomorphisms from the usual ultimate $L^{\langle-\infty\rangle}$
version of Farrell--Jones, one must also justify the change of decoration.
The vanishing conditions
\[
 \operatorname{Wh}(\pi)=0,\qquad
 \widetilde K_0(\Z\pi)=0,\qquad K_i(\Z\pi)=0\quad(i<0)
\]
provide that justification in the established reduction. They are not
consequences of the contractibility of the universal cover by elementary
homotopy theory. Nor may the orientation character be suppressed for
nonorientable manifolds. The needed implications, for torsion-free groups,
are exactly those checked in Proposition~0.3 of the cited primary paper.
\eref{115}{1}

The calculation above re-derives a known sufficient route and makes its
two-degree use explicit. It does not establish the required assembly
properties for a new class of groups. In particular, replacing them by
an unresolved global Farrell--Jones assertion makes the argument
conditional, not a proof of the catalogue's target.

\paragraph{Attempt: a constructive test on tori.}
For $f:T^n\to T^n$ a homotopy equivalence, the induced homomorphism is
some $A\in\operatorname{GL}_n(\Z)$. The linear map
\[
 h_A:\R^n/\Z^n\longrightarrow\R^n/\Z^n,
 \qquad [x]\longmapsto[Ax]
\]
is a homeomorphism. Since $T^n$ is a $K(\Z^n,1)$, maps to it are
classified up to homotopy by their induced fundamental-group homomorphisms,
up to conjugacy. The group is abelian, so $f\simeq h_A$.
This checks every homotopy class of self-equivalences of an actual torus.
It does \emph{not} show that an arbitrary closed manifold with fundamental
group $\Z^n$ is a torus; that requires the rigidity step we are studying.

Could bounded displacement of an equivariant lift yield the missing
homeomorphism directly? Already for the torus, take
\[
 F(x_1,\ldots,x_n)=
 (x_1+\sin(2\pi x_1),x_2,\ldots,x_n).
\]
It is proper, is $\Z^n$-equivariant, and stays within distance one of the
identity. It descends to a map homotopic to the identity. Nevertheless its
first-coordinate derivative is $1+2\pi\cos(2\pi x_1)$, which is positive
at $0$ and negative at $1/2$. The real first-coordinate function is not
monotone and hence is not injective. Thus $F$ is not a homeomorphism.
A coarse or bounded-distance lift construction does not by itself give
a one-to-one map; a separate controlled replacement theorem is essential.

\paragraph{Stress test.}
The source permits every dimension. Lemma~1 explicitly used the
high-dimensional topological surgery theorem, so it says nothing new about
arbitrary four-manifold groups. Even an unrestricted assembly theorem
would not allow that dimension to be silently included in this proof.
Conversely, a failure of our sufficient assembly criterion would not by
itself exhibit a nontrivial structure: actions can have stabilizers.
The calculation is deliberately one-way. \sref{115}{1}\eref{115}{1}

Contractible universal covers need not even be Euclidean; the primary
rigidity paper discusses Davis examples with that feature. Requiring an
unproved Euclidean or nonpositively curved model would exclude legitimate
inputs. Similarly, a homeomorphism after passing to finite covers is not
a descent theorem unless it conjugates the relevant deck actions.
The controlled-compactification route needs that equivariant control,
not just an abstract compactification. \eref{115}{1}\eref{115}{2}

\paragraph{Outcome.}
\textbf{Outcome: Conditional reduction.}

The notebook isolates an integral, two-degree surgery criterion, verifies
its connective-to-periodic comparison, and checks the homotopy-class
requirement in the torus test. It also gives an explicit folded equivariant
lift defeating a proposed coarse-geometric shortcut. These are
clarifications and re-derivations, not a claimed new rigidity theorem.
A complete solution still requires the requisite obstruction vanishing
for every allowed fundamental group, with a valid argument in dimension
four as well. No counterexample to the stated topological conjecture has
been constructed.

% Keep the short local bibliography together in the standalone reading copy.
\Needspace{14\baselineskip}
% END RESEARCH ADDITION 115
\subsection{Sources}
\begin{itemize}\raggedright
\item[\textnormal{[S1]}]\hypertarget{source:115:1}{} U. Hamenstadt and F. Jaeckel, Rigidity of geometric structures, Geometriae Dedicata 218, article 16 (2024), Section 2.
\url{https://link.springer.com/article/10.1007/s10711-023-00861-4}.
{\small\textit{Catalogue formulation source.}}
% BEGIN NEW REFERENCES 115
\item[\textnormal{[E1]}]\hypertarget{extra:115:1}{} Arthur Bartels and Wolfgang L\"uck,
\emph{The Borel Conjecture for hyperbolic and CAT(0)-groups},
Annals of Mathematics 175 (2012), 631--689. Theorems A--B and
Proposition 0.3, including its proof and the orientation-twisted version.
\url{https://arxiv.org/html/0901.0442v2}.
\item[\textnormal{[E2]}]\hypertarget{extra:115:2}{} Shijie Gu,
\emph{On Z-compactifiability of manifolds}, arXiv:2312.02527v3,
25 August 2026. Revised abstract and author's revision note concerning
product compatibility and additional control on universal covers.
Reported as a manuscript caution, not used as a rigidity theorem.
\url{https://arxiv.org/abs/2312.02527v3}.
% END NEW REFERENCES 115
\end{itemize}

\end{document}
